\documentclass[10pt,twocolumn]{article}
\usepackage[margin=1in]{geometry}
\usepackage{hyperref}
\usepackage{url}
\usepackage{parskip}
\usepackage{titlesec}
\usepackage{enumitem}

\title{\textbf{SEP Preliminary Report}}
\author{
  Arvin, Hanna, Jue, Yannick \\
  \small\texttt{\{Ar.Effat, Ha.Kasi, Jue.Li, Y.Radulescu\}@campus.lmu.de}
}
\date{}

\begin{document}

\maketitle

\begin{abstract}
This is our preliminary report. It covers our literature and data research,
a coarse plan of how we intend to approach the task, candidate architectures,
and our training strategy.
\end{abstract}

% ─────────────────────────────────────────────
\section{Training Data}

The goal is to find adequate, well-labeled images to train our classifier.

\subsection{Online Datasets}

Several public repositories provide pre-organized cat and dog breed datasets.
Candidate sources include:

\begin{itemize}[noitemsep]
  \item \url{https://data.mendeley.com/datasets/ng6tv57j5c/2}
  \item \url{https://www.kaggle.com/datasets/waqi786/cats-dataset}
  \item \url{https://www.kaggle.com/datasets/ma7555/cat-breeds-dataset}
\end{itemize}

A known issue with many public datasets is that they do not sort strictly by
breed, or contain images that do not belong to the target classes.
Finding a more extensive and well-curated dataset remains an open task.
As a reference point, we plan to review datasets used to train state-of-the-art
image models (e.g.\ Oxford-IIIT Pet Dataset), which address exactly our
classification problem.

\subsection{Web Scraping}

We can search for a specific breed on Google Image Search and scrape the
results. Because scraped images are noisy, a two-stage review pipeline is
needed:

\begin{enumerate}[noitemsep]
  \item \textbf{Cat/dog detector} --- filter out images that do not contain the
        target animal at all.
  \item \textbf{Breed confidence filter} --- once the model is partially
        trained on manually reviewed images, run new scraped images through it
        and accept only those above a confidence threshold for the target breed.
        These images then feed back into training.
\end{enumerate}

\subsection{AI-Generated Images}

Synthetic images produced by generative models are another data source.
Candidate generators for cats include:

\begin{itemize}[noitemsep]
  \item \url{https://theresanaiforthat.com/@izwanyassin/cute-cat-generator/}
  \item \url{https://theresanaiforthat.com/@romsike/random-adorable-cat-image-generator/}
  \item \url{https://theresanaiforthat.com/@rwsohail/cat-image-generator/}
\end{itemize}

Where available, the training datasets used by these generators may themselves
be reused. \textbf{Licensing must be verified} before incorporating any
generated or third-party data.

% ─────────────────────────────────────────────
\section{Detection Pipeline}

\subsection{Cat / Dog Detector}

We will use \textbf{YOLO} as the object detector, since it natively detects
cats and dogs. To reduce false positives and speed up inference, we will
restrict the set of detectable classes to cats and dogs only (removing all
other YOLO classes).

\subsection{Selecting the Primary Subject}

When multiple animals appear in a frame, we keep only the bounding box with
the largest area, under the assumption that the primary subject occupies the
most space in the image. If YOLO does not provide this selection directly, we
apply it as a post-processing step.

% ─────────────────────────────────────────────
\section{Image Preprocessing}

\subsection{Crop to Primary Subject}

Everything outside the selected bounding box is discarded. This removes
distracting background content and reduces the input dimensionality for the
classifier.

\subsection{Foreground Isolation (optional)}

Within the cropped region we may apply additional filtering:

\begin{itemize}[noitemsep]
  \item Remove colors that do not naturally occur in cat or dog coats.
  \item Use \textbf{Segment Anything (SAM)} to produce a pixel-precise mask of
        the animal, discarding the background entirely.
\end{itemize}

The segmentation step is optional; for speed we may skip it and pass the
cropped bounding box directly to the classifier.

% ─────────────────────────────────────────────
\section{Candidate Architectures \& Training Strategy}

\subsection{Classifier}

We plan to train a convolutional neural network (CNN) in \textbf{PyCharm /
PyTorch} on the preprocessed images. Transfer learning from a pretrained
backbone (e.g.\ ResNet or EfficientNet) is a strong candidate given our
dataset size constraints.

\subsection{Training Strategy}

\begin{enumerate}[noitemsep]
  \item \textbf{Bootstrap phase} --- train on manually reviewed images from
        the curated online datasets.
  \item \textbf{Active learning loop} --- use the partially trained model to
        score scraped and AI-generated images; add high-confidence samples to
        the training set and repeat.
  \item \textbf{Evaluation} --- hold out a balanced test set per breed;
        report top-1 and top-5 accuracy.
\end{enumerate}

\subsection{Review Workflow}

Manual review will be necessary early in the process. As the model improves,
we aim to automate the review step by using the model's own confidence score
as a proxy for image quality and label correctness.

% ─────────────────────────────────────────────
\section*{References}

\begin{itemize}[noitemsep]
  \item Parkhi et al., \textit{Cats and Dogs}, CVPR 2012
        (Oxford-IIIT Pet Dataset).
  \item Redmon et al., \textit{You Only Look Once}, CVPR 2016 (YOLO).
  \item Kirillov et al., \textit{Segment Anything}, ICCV 2023 (SAM).
\end{itemize}

\end{document}
